\section{Worked Example: Bed to Desk}

\begin{example}[A state-profile analysis rather than a fake calculation]
Consider the transition from scrolling in bed to seated reading at a desk.
\begin{enumerate}[nosep]
  \item $\Reward(s_0)$ is elevated because scrolling offers rapid novelty.
  \item $\Switch(s_0)$ is elevated because the reader must recover task
        context, locate materials, and define a first step.
  \item $\Env(s_0)$ is elevated because posture and location both stabilize
        rest-oriented behaviour.
  \item $U(s_1)$ is elevated if the study task is not pre-specified.
  \item $P$ rises when the desk is prepared, the text is opened in advance, and
        a first micro-task is written down.
\end{enumerate}
The model does not compute a true numerical barrier here. It gives a structured
explanation of why the transition feels expensive and why preparation can lower
that expense before the moment of action.
\end{example}

\begin{discussion}
This example is intentionally modest. It does not demonstrate predictive
success, parameter recovery, or empirical fit. Its purpose is narrower: to show
how the notation separates several otherwise collapsed complaints such as ``I
don't feel like it,'' ``I forgot where to start,'' and ``the room keeps pulling
me back into rest.''
\end{discussion}
